\documentclass[../../../include/open-logic-section]{subfiles}

\begin{document}

\olfileid[fr]{sth}{replacement}{strength}
\olsection{La force du remplacement}

Commençons par une observation simple sur la force du remplacement :
tant que nous ne dépassons pas $\Z$, nous ne pouvons démontrer
l'existence d'aucun ordinal de von Neumann supérieur ou égal à
$\omega + \omega$.

Voici, dans ses grandes lignes, pourquoi. En travaillant dans~$\ZF$,
considérons l'ensemble $V_{\omega+\omega}$. Il sert de domaine à un
\emph{modèle} de~$\Z$. Pour le voir, introduisons une notation pour la
\emph{relativisation} d'une formule :
\begin{defn}\ollabel{formularelativization}
	Pour tout ensemble $M$ et toute formule $\phi$, soit $\phi^M$ la formule
	obtenue en restreignant à $M$ tous les quantificateurs de $\phi$.
	Autrement dit, on remplace «~$\lexists[x]$~» par
	«~$(\lexists[x \in M])$~», et «~$\lforall[x]$~» par
	«~$(\lforall[x \in M])$~».
\end{defn}\noindent
On peut montrer que, pour chaque axiome $\phi$ de~$\Z$, on a
$\ZF \vdash \phi^{V_{\omega+\omega}}$. Mais $\omega+\omega$
n'appartient pas à $V_{\omega+\omega}$, d'après
le \olref[spine][rank]{ordsetrankalpha}. Ainsi, si $\ZF$ est cohérente, $\Z$ est compatible avec
la non-existence de $\omega+\omega$.

Voilà pourquoi nous avons dit, dans la
\olref[ordinals][replacement]{sec}, que le
\olref[ordinals][ordtype]{thmOrdinalRepresentation} ne peut être
démontré sans le remplacement. En effet, il est facile, dans~$\Z$,
de définir explicitement un bon ordre qui devrait, intuitivement,
avoir pour type d'ordre $\omega+\omega$. Nous en avons d'ailleurs donné
un exemple informel dans la \olref[ordinals][idea]{sec}, en présentant
l'ordre suivant sur les entiers naturels :
\begin{align*}
	n \lessdot m \text{ si et seulement si }&
	\text{soit }n < m\text{ et }m-n\text{ est pair,}\\
	& \text{soit $n$ est pair et $m$ est impair.}
\end{align*}
Mais si $\omega+\omega$ n'existe pas, ce bon ordre n'est isomorphe
à aucun ordinal. $\Z$ ne démontre donc \emph{pas}
le \olref[ordinals][ordtype]{thmOrdinalRepresentation}.

Inversement, le remplacement nous permet de démontrer l'existence
de $\omega+\omega$, et donc aussi celle de $V_{\omega+\omega}$.
Et ce n'est pas tout. Pour \emph{tout} bon ordre que nous pouvons
définir, le \olref[ordinals][ordtype]{thmOrdinalRepresentation}
nous assure qu'il existe un ordinal $\alpha$ isomorphe à ce bon
ordre, et donc que $V_\alpha$ existe. De manière directe, le
remplacement garantit ainsi que la hiérarchie des ensembles doit
être \emph{très haute}.

Dans les sections suivantes, puis de nouveau dans la
\olref[card-arithmetic][fix]{sec}, nous verrons mieux jusqu'à
quelle hauteur le remplacement force la hiérarchie à s'étendre.
Pour l'instant, retenons simplement que le remplacement
\emph{exige réellement} une justification !

\end{document}
